\chapter{Ретроспективное переоткрытие полной семейной связности}

Предыдущая глава обнаружила точное двумерное ядро директорного символа.
Теперь необходимо проверить, не является ли оно указанием на уже
возникавшую в проекте структуру, а не требованием вводить произвольный
квадратичный член. Ретроспектива действительно выделяет один строго
ограниченный кандидат: вещественную семейную $SO(3)$-связность.

\section{Что из прошлых результатов сохраняет силу}

Проверены близкие ветви Томов IV и V. Их выводы разделяются на две группы.

Положительно используются следующие результаты.
\begin{enumerate}
 \item Семейный носитель уже содержит каноническое вещественное
 трёхмерное представление $SO(3)$.
 \item В локальной трубчатой модели связность этого представления давала
 точное одномерное ядро и положительную щель семейных мод.
 \item Поле порядка $Q$ преобразуется в пятимерном представлении
 симметричных бесследовых тензоров того же $SO(3)$.
 \item Составная связность $A_Q=[Q,dQ]/\Delta^2$ точно восстанавливает две
 косетные компоненты $SO(3)/O(2)$, но пропускает непрерывный генератор
 стабилизатора.
\end{enumerate}

Одновременно сохраняются прежние запреты.
\begin{enumerate}
 \item Полные матричные контейнеры $M_{35}$ и $M_{300}$ нельзя объявлять
 координатными калибровочными алгебрами.
 \item След сам по себе не выбирает пространственную связность, её
 вложение и абсолютную нормировку.
 \item Прямоугольный нечётный носитель $20\times15$ не даёт гладкий
 спинорный оператор Каллиаса с единичным индексом.
 \item Обычные внутренние одноформы Стандартной модели не видят семейный
 оператор.
\end{enumerate}

Последний пункт о спинорном носителе не запрещает бозонную связность,
действующую на уже существующих вещественных представлениях $3$ и $5$.
Именно это различение ранее не было нужно, но стало решающим после
обнаружения директорного ядра.

\section{Точная орбитальная карта вакуума}

Пусть
\begin{equation}
 Q_0=\Delta\left(P_0-\frac{I_3}{3}\right),
 \qquad P_0=\operatorname{diag}(1,0,0).
\end{equation}
Линеаризованное действие алгебры $so(3)$ задаётся картой
\begin{equation}
 L_{Q_0}:so(3)\longrightarrow\operatorname{Sym}_0(3),
 \qquad \omega\longmapsto[\omega,Q_0].
 \label{eq:v6-full-gauge-orbit-map}
\end{equation}
В ортонормированных базисах её ненулевые сингулярные числа равны
\begin{equation}
 \Spec_{\rm sing}(L_{Q_0})=\{\Delta,\Delta,0\}.
\end{equation}
Следовательно,
\begin{equation}
 \rank L_{Q_0}=2,
 \qquad \dim\ker L_{Q_0}=1.
\end{equation}
Образ карты в точности совпадает с двумя директорными направлениями,
которые образовали ядро предыдущего пятикомпонентного символа. Ядро карты
есть непрерывная алгебра стабилизатора $so(2)\subset O(2)$.

Квадрат отображения на калибровочной алгебре имеет спектр
\begin{equation}
 \Spec(L_{Q_0}^{*}L_{Q_0})
 =\{0,\Delta^2,\Delta^2\}
 =\{0,0.7538578897\ldots,0.7538578897\ldots\}.
 \label{eq:v6-full-gauge-mass-map}
\end{equation}
Это стандартная схема нарушения $SO(3)\to O(2)$: две компоненты связности
видят вакуум, а одна непрерывная стабилизаторная компонента остаётся
безмассовой на однородном фоне.

\section{Минимальное полное завершение}

Введём независимую пространственную связность
\begin{equation}
 B_i\in so(3),
 \qquad
 D_iQ=\partial_iQ+[B_i,Q],
\end{equation}
с законом локального преобразования
\begin{equation}
 \delta_\omega Q=[\omega,Q],
 \qquad
 \delta_\omega B_i=[\omega,B_i]-\partial_i\omega.
 \label{eq:v6-full-gauge-law}
\end{equation}
На постоянном вакууме линейная чисто калибровочная вариация
\begin{equation}
 \delta Q=[\omega,Q_0],
 \qquad
 \delta B_i=-\partial_i\omega
\end{equation}
даёт
\begin{equation}
 \delta(D_iQ)=0.
\end{equation}
Машинный остаток этого тождества для общего волнового вектора равен
$1.4\cdot10^{-17}$.

Тем самым точные одномерные директорные семейства предыдущей главы
получают новое чтение. После добавления независимой полной связности они
являются локальными калибровочными образами постоянного вакуума с плоской
связностью, а не физически различными вакуумами.

\section{Калибровочно фиксированный символ}

Рассмотрим минимальный квадратичный функционал
\begin{equation}
 \mathcal E^{(2)}=
 \sum_i\|\partial_i\delta Q+[\delta B_i,Q_0]\|^2
 +\frac12\sum_{i,j}\|\partial_i\delta B_j-\partial_j\delta B_i\|^2
 +\frac12\langle\delta Q,H_V\delta Q\rangle,
 \label{eq:v6-full-gauge-quadratic-energy}
\end{equation}
где $H_V$ есть уже вычисленный гессиан канонического потенциала. Для
проверки символа добавлено условие типа 'т Хоофта
\begin{equation}
 \mathcal G=\partial_i\delta B_i-L_{Q_0}^{*}\delta Q.
 \label{eq:v6-full-gauge-fixing}
\end{equation}

Совместный символ действует на четырнадцати компонентах: пяти компонентах
$\delta Q$ и девяти компонентах $\delta B_i$. При волновых числах
$k=1/4,1/2,1,3$ его ранг равен четырнадцати. Наименьшие собственные
значения равны соответственно
\begin{equation}
 \frac1{16},\qquad\frac14,\qquad1,\qquad9.
\end{equation}
То есть они совпадают с $k^2$ стабилизаторного максвелловского сектора.
Директорное ядро устранено калибровочным факторированием, а не
искусственной массой.

\begin{theorem}[Кинематическое переоткрытие полной связности]
Полная независимая $SO(3)$-связность является минимальным локальным
завершением, точно соответствующим обнаруженному двумерному директорному
ядру. После стандартной калибровочной фиксации вакуумный символ имеет
полный ранг при каждом проверенном $k\ne0$. Для этого бозонного результата
не требуется спинорный дублет или ремонт прямоугольного носителя
$20\times15$.
\end{theorem}

Этот результат согласуется с обычным механизмом Браута--Энглерта--Хиггса
и с теориями неабелевых дефектов с ориентационными модами; см. DOI
\path{10.1016/0550-3213(74)90486-6}, DOI
\path{10.1070/PU1974v017n03ABEH004089},
\path{arXiv:hep-th/0307287} и \path{arXiv:1509.04061}.

\section{Почему физическое замыкание ещё не достигнуто}

Проверка устанавливает необходимость и кинематическую достаточность типа
поля, но не выводит его из родителя. Независимая связность имеет
собственную кривизну и потому означает физический калибровочный сектор:
две его компоненты становятся массивными на упорядоченном фоне, а
непрерывная стабилизаторная компонента остаётся безмассовой до дальнейшего
нарушения или ограничения.

Нельзя считать этот сектор уже полученным только потому, что:
\begin{enumerate}
 \item составная косетная часть $A_Q$ имеет правильную формулу;
 \item локальная семейная связность ранее использовалась в трубчатой
 модели;
 \item след $M_{35}$ способен нормировать действие после задания
 вложения.
\end{enumerate}
Требуется единое отождествление пространственной связности $B$ с прежней
семейной связностью, совместимое с глобальным носителем, мерой и
наблюдаемым спектром. Иначе мы просто добавили три новых калибровочных
поля после обнаружения проблемы.

\begin{nogocriterion}[Происхождение связности остаётся открытым]
Кинематическая эллиптичность не считается выводом бозонного сектора, пока
из уже принятой архитектуры не получены: локальный закон семейной
$SO(3)$-симметрии, кривизный член с общей нормировкой, глобальное
продолжение через дефект и судьба непрерывной стабилизаторной компоненты.
\end{nogocriterion}

\section{Вердикт}

\begin{protocol}
\begin{enumerate}
 \item ретроспективная сверка с Томами IV--V: \textbf{выполнена};
 \item образ $so(3)$ на вакууме: \textbf{ровно две директорные моды};
 \item ядро орбитальной карты: \textbf{один генератор стабилизатора};
 \item линейная калибровочная ковариантность: \textbf{проверена};
 \item ранг фиксированного символа при $k\ne0$: \textbf{$14$ из $14$};
 \item точные плоские директорные семейства: \textbf{чистая калибровка};
 \item спинорный дублет для бозонной эллиптичности: \textbf{не требуется};
 \item прежний запрет $20\times15$: \textbf{сохраняется для фермионов};
 \item полные $M_{35}$ или $M_{300}$ как калибровочные алгебры:
 \textbf{не возвращены};
 \item родительское происхождение независимой связности: \textbf{открыто};
 \item рождение материи: \textbf{не закрыто};
 \item следующий гейт: \textbf{отождествление полной связности с уже
 существующей семейной связностью родителя}.
\end{enumerate}
\end{protocol}

Машинный сертификат сохранён в
\path{s2t_v6_bosonic_defect_full_gauge_completion_reopening_gate_results.json}.